% method_theory.tex — Method and Theoretical Analysis sections
% Include from paper.tex via % method_theory.tex — Method and Theoretical Analysis sections
% Include from paper.tex via % method_theory.tex — Method and Theoretical Analysis sections
% Include from paper.tex via % method_theory.tex — Method and Theoretical Analysis sections
% Include from paper.tex via \input{sections/method_theory}

% ===========================================================================
\section{Method}
\label{sec:method}

% ---------------------------------------------------------------------------
\subsection{Problem Setup}
\label{sec:method:setup}

Let $\calX$ and $\calY$ denote spaces of functions defined on a bounded domain
$\Omega \subset \R^d$.  A neural operator $G_\theta : \calX \to \calY$,
parameterised by $\theta$, maps an input field (e.g., initial condition,
forcing term, or coefficient field) to a solution field of a partial
differential equation.  We train $G_\theta$ on a dataset
$\{(x_i, y_i)\}_{i=1}^{N}$ of input--solution pairs generated by a
numerical PDE solver.

After training, our goal is to construct \emph{prediction bands}
$\calC_\alpha(x) \subseteq \calY$ such that, for a new exchangeable test
point $(X_{n+1}, Y_{n+1})$,
\begin{equation}
  \label{eq:coverage}
  \Pr\!\bigl(Y_{n+1} \in \calC_\alpha(X_{n+1})\bigr) \;\geq\; 1 - \alpha,
\end{equation}
where $\alpha \in (0,1)$ is the user-specified miscoverage level.
Crucially, this guarantee must hold in finite samples and without
distributional assumptions on the data-generating process, beyond
exchangeability.

% ---------------------------------------------------------------------------
\subsection{Spectral Nonconformity Scores}
\label{sec:method:scores}

Given a prediction $\hat{y} = G_\theta(x)$ and the true solution $y$,
define the \emph{error field}
\begin{equation}
  e(x) \;=\; y - G_\theta(x).
\end{equation}
We analyse $e$ in the frequency domain via the discrete Fourier transform
(DFT):
\begin{equation}
  \hat{e}(k) \;=\; \mathrm{FFT}(e)(k), \qquad k \in \mathbb{Z}^d.
\end{equation}

\begin{definition}[Weighted Spectral Score]
\label{def:spectral-score}
For a weight function $w : \mathbb{Z}^d \to [0, \infty)$, the
\emph{weighted spectral nonconformity score} is
\begin{equation}
  \label{eq:spectral-score}
  \sSpec^{w}(G_\theta(x),\, y)
  \;=\; \sum_{k} w_k \,\bigl|\hat{e}(k)\bigr|^2.
\end{equation}
\end{definition}

Different choices of~$w$ recover well-known norms as special cases:
\begin{itemize}
  \item \textbf{$L^2$ score.}  Setting $w_k = 1$ for all~$k$ yields, by
    Parseval's theorem,
    $\sSpec^{1}(G_\theta(x), y) = \|y - G_\theta(x)\|_{L^2}^2$.
  \item \textbf{Sobolev $H^s$ score.}  Setting
    $w_k = (1 + |k|^2)^{s}$ for $s > 0$ yields the squared $H^s$
    Sobolev semi-norm of the error, which penalises high-frequency
    discrepancies more heavily.
  \item \textbf{Adaptive (learned) score.}  Setting $w_k$ to
    data-driven values (Section~\ref{sec:method:learned}) allows the
    score to concentrate on the frequency bands where the neural operator
    errs most, producing tighter prediction bands.
\end{itemize}

% ---------------------------------------------------------------------------
\subsection{Split Conformal Prediction with Spectral Scores}
\label{sec:method:conformal}

We apply the standard split conformal protocol
\citep{vovk2005algorithmic,lei2018distribution} with the spectral score
\eqref{eq:spectral-score} as the nonconformity measure.

Given a calibration set $\mathcal{D}_{\mathrm{cal}} = \{(x_i, y_i)\}_{i=1}^{n}$
that is exchangeable with the test point, we compute scores
$s_i = \sSpec^{w}(G_\theta(x_i), y_i)$ for $i = 1, \dots, n$ and set
\begin{equation}
  \label{eq:quantile}
  \qhat \;=\; \mathrm{Quantile}\!\Bigl(
    s_1, \dots, s_n;\;
    \frac{\lceil (1-\alpha)(n+1) \rceil}{n}
  \Bigr).
\end{equation}
The prediction band for a new input $x$ is then
\begin{equation}
  \label{eq:pred-band}
  \calC_\alpha(x)
  \;=\;
  \bigl\{\, y \in \calY : \sSpec^{w}(G_\theta(x),\, y) \leq \qhat \,\bigr\}.
\end{equation}
Because the spectral score is a measurable function of the residual, the
standard exchangeability-based coverage guarantee
(Theorem~\ref{thm:coverage}) applies immediately.

% ---------------------------------------------------------------------------
\subsection{Learned Spectral Weights}
\label{sec:method:learned}

The choice of weight vector~$w$ does not affect the \emph{validity} of
coverage (any measurable score preserves the guarantee), but it
strongly influences \emph{efficiency}: the average width of the
prediction band.  We therefore propose learning~$w$ from data.

\paragraph{Three-way data split.}
To avoid statistical double-dipping, we partition the held-out data into
three disjoint subsets:
\begin{enumerate}
  \item $\mathcal{D}_w$: used to learn the weight vector~$w$;
  \item $\mathcal{D}_{\mathrm{cal}}$: used to compute the conformal
    quantile $\qhat$ with the learned weights;
  \item $\mathcal{D}_{\mathrm{test}}$: used for evaluation only.
\end{enumerate}
The three-way split is necessary because using a single calibration set both
to optimise~$w$ and to compute~$\qhat$ would violate exchangeability
between the calibration scores and the test score, invalidating the
coverage guarantee.

\paragraph{Objective.}
On $\mathcal{D}_w$, we compute per-sample spectral error power spectra
and seek weights that minimise the \emph{variance} of the resulting scores:
\begin{equation}
  \label{eq:weight-obj}
  w^{\star}
  \;=\;
  \operatorname*{arg\,min}_{w \geq 0}
  \;\Var_{(x,y) \in \mathcal{D}_w}\!\bigl[\sSpec^{w}(G_\theta(x), y)\bigr].
\end{equation}
Minimising score variance concentrates the score distribution, which in turn
yields a smaller conformal quantile $\qhat$ and therefore tighter prediction
bands for a given coverage level~$1-\alpha$.

\paragraph{Binned parameterisation.}
To keep the optimisation low-dimensional and well-conditioned, we do not
learn one weight per frequency~$k$.  Instead, we partition the frequency
indices into $B = 16$ logarithmically spaced bins and assign a single
scalar weight to all frequencies within each bin:
$w_k = w_{b(k)}$ where $b(k) \in \{1, \dots, B\}$.
The resulting 16-dimensional optimisation is solved by L-BFGS-B with
non-negativity constraints.

% ---------------------------------------------------------------------------
\subsection{Algorithm}
\label{sec:method:algorithm}

Algorithm~\ref{alg:spectral-conformal} summarises the complete pipeline.

\begin{algorithm}[t]
\DontPrintSemicolon
\SetAlgoLined
\KwIn{Training data $\mathcal{D}_{\mathrm{train}}$; held-out data
  $\mathcal{D}_{\mathrm{held}}$; miscoverage level $\alpha$;
  number of frequency bins $B$}
\KwOut{Prediction band function $\calC_\alpha(\cdot)$}

\BlankLine
\tcp{Step 1: Train neural operator}
Train $G_\theta$ on $\mathcal{D}_{\mathrm{train}}$\;

\BlankLine
\tcp{Step 2: Three-way split of held-out data}
$\mathcal{D}_w,\, \mathcal{D}_{\mathrm{cal}},\, \mathcal{D}_{\mathrm{test}}
  \leftarrow \mathrm{Split}(\mathcal{D}_{\mathrm{held}})$\;

\BlankLine
\tcp{Step 3: Learn spectral weights}
\For{$(x_i, y_i) \in \mathcal{D}_w$}{
  $e_i \leftarrow y_i - G_\theta(x_i)$\;
  $\hat{e}_i \leftarrow \mathrm{FFT}(e_i)$\;
  Accumulate $|\hat{e}_i(k)|^2$ per frequency bin\;
}
$w^{\star} \leftarrow \operatorname*{arg\,min}_{w \geq 0}
  \Var\!\bigl[\sSpec^{w}\bigr]$ on $\mathcal{D}_w$
  \tcp*{Eq.~\eqref{eq:weight-obj}}

\BlankLine
\tcp{Step 4: Calibrate}
\For{$(x_i, y_i) \in \mathcal{D}_{\mathrm{cal}}$}{
  $s_i \leftarrow \sSpec^{w^{\star}}(G_\theta(x_i),\, y_i)$
    \tcp*{Eq.~\eqref{eq:spectral-score}}
}
$\qhat \leftarrow \mathrm{Quantile}\!\bigl(
  s_1, \dots, s_n;\;
  \lceil(1-\alpha)(n+1)\rceil / n\bigr)$
  \tcp*{Eq.~\eqref{eq:quantile}}

\BlankLine
\tcp{Step 5: Predict}
\For{new input $x$}{
  $\calC_\alpha(x) \leftarrow
    \{y : \sSpec^{w^{\star}}(G_\theta(x), y) \leq \qhat\}$\;
}
\Return{$\calC_\alpha$}\;

\caption{Spectral Conformal Prediction for Neural Operators}
\label{alg:spectral-conformal}
\end{algorithm}


% ===========================================================================
\section{Theoretical Analysis}
\label{sec:theory}

We establish three results: (i)~the spectral score with uniform weights
recovers the standard $L^2$ score, (ii)~the coverage guarantee of split
conformal prediction extends to all spectral scores, and (iii)~spectral
weighting can provably tighten prediction bands when errors exhibit
frequency concentration.

% ---------------------------------------------------------------------------
\begin{theorem}[Parseval Consistency]
\label{thm:parseval}
For uniform weights $w_k = 1$ for all~$k$, the weighted spectral score
equals the squared $L^2$ norm of the residual:
\begin{equation}
  \sSpec^{1}(G_\theta(x),\, y)
  \;=\;
  \|y - G_\theta(x)\|_{L^2(\Omega)}^2.
\end{equation}
\end{theorem}

\begin{proof}
By Parseval's theorem for the discrete Fourier transform,
\begin{equation}
  \sum_{k} \bigl|\hat{e}(k)\bigr|^2
  \;=\;
  \sum_{j} \bigl|e(j)\bigr|^2
  \;=\;
  \|e\|_{L^2}^2,
\end{equation}
where $e = y - G_\theta(x)$.  Substituting into
Definition~\ref{def:spectral-score} with $w_k = 1$ yields the result.
\end{proof}

% ---------------------------------------------------------------------------
\begin{theorem}[Coverage Guarantee]
\label{thm:coverage}
Let $s : \calX \times \calY \to \R$ be any measurable nonconformity score
function, and let $(X_1, Y_1), \dots, (X_n, Y_n), (X_{n+1}, Y_{n+1})$
be exchangeable.  Define $s_i = s(X_i, Y_i)$ and
\begin{equation}
  \qhat = \inf\!\Bigl\{q : \frac{|\{i \leq n : s_i \leq q\}|}{n}
    \geq \frac{\lceil (1-\alpha)(n+1) \rceil}{n}\Bigr\}.
\end{equation}
Then the prediction set $\calC_\alpha(x) = \{y : s(x,y) \leq \qhat\}$
satisfies
\begin{equation}
  \Pr\!\bigl(Y_{n+1} \in \calC_\alpha(X_{n+1})\bigr)
  \;\geq\;
  1 - \alpha.
\end{equation}
\end{theorem}

\begin{proof}
The proof follows directly from the standard split conformal guarantee
\citep{vovk2005algorithmic,lei2018distribution}.  The key observation is
that exchangeability of the data pairs $(X_i, Y_i)$ implies
exchangeability of the scores $s_i$, since $s$ is a fixed measurable
function applied identically to each pair.  In particular, the spectral
score $\sSpec^{w}$ is measurable because the FFT is a linear (hence
measurable) map and the weighted sum of squared moduli is a continuous
function of the Fourier coefficients.  The rank of $s_{n+1}$ among
$\{s_1, \dots, s_n, s_{n+1}\}$ is therefore uniformly distributed over
$\{1, \dots, n+1\}$, from which the coverage bound follows by the
standard quantile argument.
\end{proof}

% ---------------------------------------------------------------------------
\begin{proposition}[Band Width Reduction via Spectral Concentration]
\label{prop:width}
Let $\hat{e}_1, \dots, \hat{e}_n$ be the Fourier-domain error fields on
the calibration set.  Suppose that the spectral energy of the errors is
concentrated in a subset of frequency indices $K_0 \subset \mathbb{Z}^d$:
\begin{equation}
  \frac{\E\!\bigl[\sum_{k \in K_0} |\hat{e}(k)|^2\bigr]}
       {\E\!\bigl[\sum_{k} |\hat{e}(k)|^2\bigr]}
  \;\geq\; 1 - \delta,
  \qquad \delta \ll 1.
\end{equation}
Let $w^{\star}$ be the variance-minimising weights
\eqref{eq:weight-obj} and $w^{1}_k = 1$ the uniform weights.
Then the conformal quantile satisfies
$\qhat(w^{\star}) \leq \qhat(w^{1})$, with equality only when the
spectral energy is uniformly distributed across all frequencies.
Consequently, the prediction band under $w^{\star}$ is no wider than
under uniform weights, at the same coverage level~$1-\alpha$.
\end{proposition}

\begin{proof}[Proof sketch]
When errors concentrate in $K_0$, the spectral power outside $K_0$
contributes primarily noise variance to the scores under uniform weights.
By assigning larger weights to $K_0$ and down-weighting the remaining
frequencies, the optimised score $\sSpec^{w^{\star}}$ has strictly lower
variance than $\sSpec^{1}$ (the overall scale of scores is absorbed by the
quantile computation).

More precisely, write $s^{w} = w^{\!\top} p$ where
$p_k = |\hat{e}(k)|^2$ is the spectral power vector.  Then
$\Var[s^{w}] = w^{\!\top} \Cov[p]\, w$.  Under spectral concentration,
$\Cov[p]$ has eigenvalues that are large only in the directions
corresponding to~$K_0$; aligning~$w$ with these directions while
suppressing the noisy tail reduces $\Var[s^{w}]$.  A lower-variance
score distribution yields a tighter gap between its median and its
$(1-\alpha)$-quantile, so $\qhat(w^{\star}) \leq \qhat(w^{1})$, and the
prediction band $\{y : \sSpec^{w^{\star}} \leq \qhat(w^{\star})\}$ is
correspondingly tighter.  The coverage guarantee is preserved because
the weights are fixed before calibration
(Section~\ref{sec:method:learned}).
\end{proof}

\begin{remark}
The three-way data split is essential for
Proposition~\ref{prop:width} to coexist with Theorem~\ref{thm:coverage}.
If the same data were used both to learn~$w$ and to compute~$\qhat$,
the calibration scores would no longer be exchangeable with the test
score, and the coverage guarantee would be invalidated.
\end{remark}


% ===========================================================================
\section{Method}
\label{sec:method}

% ---------------------------------------------------------------------------
\subsection{Problem Setup}
\label{sec:method:setup}

Let $\calX$ and $\calY$ denote spaces of functions defined on a bounded domain
$\Omega \subset \R^d$.  A neural operator $G_\theta : \calX \to \calY$,
parameterised by $\theta$, maps an input field (e.g., initial condition,
forcing term, or coefficient field) to a solution field of a partial
differential equation.  We train $G_\theta$ on a dataset
$\{(x_i, y_i)\}_{i=1}^{N}$ of input--solution pairs generated by a
numerical PDE solver.

After training, our goal is to construct \emph{prediction bands}
$\calC_\alpha(x) \subseteq \calY$ such that, for a new exchangeable test
point $(X_{n+1}, Y_{n+1})$,
\begin{equation}
  \label{eq:coverage}
  \Pr\!\bigl(Y_{n+1} \in \calC_\alpha(X_{n+1})\bigr) \;\geq\; 1 - \alpha,
\end{equation}
where $\alpha \in (0,1)$ is the user-specified miscoverage level.
Crucially, this guarantee must hold in finite samples and without
distributional assumptions on the data-generating process, beyond
exchangeability.

% ---------------------------------------------------------------------------
\subsection{Spectral Nonconformity Scores}
\label{sec:method:scores}

Given a prediction $\hat{y} = G_\theta(x)$ and the true solution $y$,
define the \emph{error field}
\begin{equation}
  e(x) \;=\; y - G_\theta(x).
\end{equation}
We analyse $e$ in the frequency domain via the discrete Fourier transform
(DFT):
\begin{equation}
  \hat{e}(k) \;=\; \mathrm{FFT}(e)(k), \qquad k \in \mathbb{Z}^d.
\end{equation}

\begin{definition}[Weighted Spectral Score]
\label{def:spectral-score}
For a weight function $w : \mathbb{Z}^d \to [0, \infty)$, the
\emph{weighted spectral nonconformity score} is
\begin{equation}
  \label{eq:spectral-score}
  \sSpec^{w}(G_\theta(x),\, y)
  \;=\; \sum_{k} w_k \,\bigl|\hat{e}(k)\bigr|^2.
\end{equation}
\end{definition}

Different choices of~$w$ recover well-known norms as special cases:
\begin{itemize}
  \item \textbf{$L^2$ score.}  Setting $w_k = 1$ for all~$k$ yields, by
    Parseval's theorem,
    $\sSpec^{1}(G_\theta(x), y) = \|y - G_\theta(x)\|_{L^2}^2$.
  \item \textbf{Sobolev $H^s$ score.}  Setting
    $w_k = (1 + |k|^2)^{s}$ for $s > 0$ yields the squared $H^s$
    Sobolev semi-norm of the error, which penalises high-frequency
    discrepancies more heavily.
  \item \textbf{Adaptive (learned) score.}  Setting $w_k$ to
    data-driven values (Section~\ref{sec:method:learned}) allows the
    score to concentrate on the frequency bands where the neural operator
    errs most, producing tighter prediction bands.
\end{itemize}

% ---------------------------------------------------------------------------
\subsection{Split Conformal Prediction with Spectral Scores}
\label{sec:method:conformal}

We apply the standard split conformal protocol
\citep{vovk2005algorithmic,lei2018distribution} with the spectral score
\eqref{eq:spectral-score} as the nonconformity measure.

Given a calibration set $\mathcal{D}_{\mathrm{cal}} = \{(x_i, y_i)\}_{i=1}^{n}$
that is exchangeable with the test point, we compute scores
$s_i = \sSpec^{w}(G_\theta(x_i), y_i)$ for $i = 1, \dots, n$ and set
\begin{equation}
  \label{eq:quantile}
  \qhat \;=\; \mathrm{Quantile}\!\Bigl(
    s_1, \dots, s_n;\;
    \frac{\lceil (1-\alpha)(n+1) \rceil}{n}
  \Bigr).
\end{equation}
The prediction band for a new input $x$ is then
\begin{equation}
  \label{eq:pred-band}
  \calC_\alpha(x)
  \;=\;
  \bigl\{\, y \in \calY : \sSpec^{w}(G_\theta(x),\, y) \leq \qhat \,\bigr\}.
\end{equation}
Because the spectral score is a measurable function of the residual, the
standard exchangeability-based coverage guarantee
(Theorem~\ref{thm:coverage}) applies immediately.

% ---------------------------------------------------------------------------
\subsection{Learned Spectral Weights}
\label{sec:method:learned}

The choice of weight vector~$w$ does not affect the \emph{validity} of
coverage (any measurable score preserves the guarantee), but it
strongly influences \emph{efficiency}: the average width of the
prediction band.  We therefore propose learning~$w$ from data.

\paragraph{Three-way data split.}
To avoid statistical double-dipping, we partition the held-out data into
three disjoint subsets:
\begin{enumerate}
  \item $\mathcal{D}_w$: used to learn the weight vector~$w$;
  \item $\mathcal{D}_{\mathrm{cal}}$: used to compute the conformal
    quantile $\qhat$ with the learned weights;
  \item $\mathcal{D}_{\mathrm{test}}$: used for evaluation only.
\end{enumerate}
The three-way split is necessary because using a single calibration set both
to optimise~$w$ and to compute~$\qhat$ would violate exchangeability
between the calibration scores and the test score, invalidating the
coverage guarantee.

\paragraph{Objective.}
On $\mathcal{D}_w$, we compute per-sample spectral error power spectra
and seek weights that minimise the \emph{variance} of the resulting scores:
\begin{equation}
  \label{eq:weight-obj}
  w^{\star}
  \;=\;
  \operatorname*{arg\,min}_{w \geq 0}
  \;\Var_{(x,y) \in \mathcal{D}_w}\!\bigl[\sSpec^{w}(G_\theta(x), y)\bigr].
\end{equation}
Minimising score variance concentrates the score distribution, which in turn
yields a smaller conformal quantile $\qhat$ and therefore tighter prediction
bands for a given coverage level~$1-\alpha$.

\paragraph{Binned parameterisation.}
To keep the optimisation low-dimensional and well-conditioned, we do not
learn one weight per frequency~$k$.  Instead, we partition the frequency
indices into $B = 16$ logarithmically spaced bins and assign a single
scalar weight to all frequencies within each bin:
$w_k = w_{b(k)}$ where $b(k) \in \{1, \dots, B\}$.
The resulting 16-dimensional optimisation is solved by L-BFGS-B with
non-negativity constraints.

% ---------------------------------------------------------------------------
\subsection{Algorithm}
\label{sec:method:algorithm}

Algorithm~\ref{alg:spectral-conformal} summarises the complete pipeline.

\begin{algorithm}[t]
\DontPrintSemicolon
\SetAlgoLined
\KwIn{Training data $\mathcal{D}_{\mathrm{train}}$; held-out data
  $\mathcal{D}_{\mathrm{held}}$; miscoverage level $\alpha$;
  number of frequency bins $B$}
\KwOut{Prediction band function $\calC_\alpha(\cdot)$}

\BlankLine
\tcp{Step 1: Train neural operator}
Train $G_\theta$ on $\mathcal{D}_{\mathrm{train}}$\;

\BlankLine
\tcp{Step 2: Three-way split of held-out data}
$\mathcal{D}_w,\, \mathcal{D}_{\mathrm{cal}},\, \mathcal{D}_{\mathrm{test}}
  \leftarrow \mathrm{Split}(\mathcal{D}_{\mathrm{held}})$\;

\BlankLine
\tcp{Step 3: Learn spectral weights}
\For{$(x_i, y_i) \in \mathcal{D}_w$}{
  $e_i \leftarrow y_i - G_\theta(x_i)$\;
  $\hat{e}_i \leftarrow \mathrm{FFT}(e_i)$\;
  Accumulate $|\hat{e}_i(k)|^2$ per frequency bin\;
}
$w^{\star} \leftarrow \operatorname*{arg\,min}_{w \geq 0}
  \Var\!\bigl[\sSpec^{w}\bigr]$ on $\mathcal{D}_w$
  \tcp*{Eq.~\eqref{eq:weight-obj}}

\BlankLine
\tcp{Step 4: Calibrate}
\For{$(x_i, y_i) \in \mathcal{D}_{\mathrm{cal}}$}{
  $s_i \leftarrow \sSpec^{w^{\star}}(G_\theta(x_i),\, y_i)$
    \tcp*{Eq.~\eqref{eq:spectral-score}}
}
$\qhat \leftarrow \mathrm{Quantile}\!\bigl(
  s_1, \dots, s_n;\;
  \lceil(1-\alpha)(n+1)\rceil / n\bigr)$
  \tcp*{Eq.~\eqref{eq:quantile}}

\BlankLine
\tcp{Step 5: Predict}
\For{new input $x$}{
  $\calC_\alpha(x) \leftarrow
    \{y : \sSpec^{w^{\star}}(G_\theta(x), y) \leq \qhat\}$\;
}
\Return{$\calC_\alpha$}\;

\caption{Spectral Conformal Prediction for Neural Operators}
\label{alg:spectral-conformal}
\end{algorithm}


% ===========================================================================
\section{Theoretical Analysis}
\label{sec:theory}

We establish three results: (i)~the spectral score with uniform weights
recovers the standard $L^2$ score, (ii)~the coverage guarantee of split
conformal prediction extends to all spectral scores, and (iii)~spectral
weighting can provably tighten prediction bands when errors exhibit
frequency concentration.

% ---------------------------------------------------------------------------
\begin{theorem}[Parseval Consistency]
\label{thm:parseval}
For uniform weights $w_k = 1$ for all~$k$, the weighted spectral score
equals the squared $L^2$ norm of the residual:
\begin{equation}
  \sSpec^{1}(G_\theta(x),\, y)
  \;=\;
  \|y - G_\theta(x)\|_{L^2(\Omega)}^2.
\end{equation}
\end{theorem}

\begin{proof}
By Parseval's theorem for the discrete Fourier transform,
\begin{equation}
  \sum_{k} \bigl|\hat{e}(k)\bigr|^2
  \;=\;
  \sum_{j} \bigl|e(j)\bigr|^2
  \;=\;
  \|e\|_{L^2}^2,
\end{equation}
where $e = y - G_\theta(x)$.  Substituting into
Definition~\ref{def:spectral-score} with $w_k = 1$ yields the result.
\end{proof}

% ---------------------------------------------------------------------------
\begin{theorem}[Coverage Guarantee]
\label{thm:coverage}
Let $s : \calX \times \calY \to \R$ be any measurable nonconformity score
function, and let $(X_1, Y_1), \dots, (X_n, Y_n), (X_{n+1}, Y_{n+1})$
be exchangeable.  Define $s_i = s(X_i, Y_i)$ and
\begin{equation}
  \qhat = \inf\!\Bigl\{q : \frac{|\{i \leq n : s_i \leq q\}|}{n}
    \geq \frac{\lceil (1-\alpha)(n+1) \rceil}{n}\Bigr\}.
\end{equation}
Then the prediction set $\calC_\alpha(x) = \{y : s(x,y) \leq \qhat\}$
satisfies
\begin{equation}
  \Pr\!\bigl(Y_{n+1} \in \calC_\alpha(X_{n+1})\bigr)
  \;\geq\;
  1 - \alpha.
\end{equation}
\end{theorem}

\begin{proof}
The proof follows directly from the standard split conformal guarantee
\citep{vovk2005algorithmic,lei2018distribution}.  The key observation is
that exchangeability of the data pairs $(X_i, Y_i)$ implies
exchangeability of the scores $s_i$, since $s$ is a fixed measurable
function applied identically to each pair.  In particular, the spectral
score $\sSpec^{w}$ is measurable because the FFT is a linear (hence
measurable) map and the weighted sum of squared moduli is a continuous
function of the Fourier coefficients.  The rank of $s_{n+1}$ among
$\{s_1, \dots, s_n, s_{n+1}\}$ is therefore uniformly distributed over
$\{1, \dots, n+1\}$, from which the coverage bound follows by the
standard quantile argument.
\end{proof}

% ---------------------------------------------------------------------------
\begin{proposition}[Band Width Reduction via Spectral Concentration]
\label{prop:width}
Let $\hat{e}_1, \dots, \hat{e}_n$ be the Fourier-domain error fields on
the calibration set.  Suppose that the spectral energy of the errors is
concentrated in a subset of frequency indices $K_0 \subset \mathbb{Z}^d$:
\begin{equation}
  \frac{\E\!\bigl[\sum_{k \in K_0} |\hat{e}(k)|^2\bigr]}
       {\E\!\bigl[\sum_{k} |\hat{e}(k)|^2\bigr]}
  \;\geq\; 1 - \delta,
  \qquad \delta \ll 1.
\end{equation}
Let $w^{\star}$ be the variance-minimising weights
\eqref{eq:weight-obj} and $w^{1}_k = 1$ the uniform weights.
Then the conformal quantile satisfies
$\qhat(w^{\star}) \leq \qhat(w^{1})$, with equality only when the
spectral energy is uniformly distributed across all frequencies.
Consequently, the prediction band under $w^{\star}$ is no wider than
under uniform weights, at the same coverage level~$1-\alpha$.
\end{proposition}

\begin{proof}[Proof sketch]
When errors concentrate in $K_0$, the spectral power outside $K_0$
contributes primarily noise variance to the scores under uniform weights.
By assigning larger weights to $K_0$ and down-weighting the remaining
frequencies, the optimised score $\sSpec^{w^{\star}}$ has strictly lower
variance than $\sSpec^{1}$ (the overall scale of scores is absorbed by the
quantile computation).

More precisely, write $s^{w} = w^{\!\top} p$ where
$p_k = |\hat{e}(k)|^2$ is the spectral power vector.  Then
$\Var[s^{w}] = w^{\!\top} \Cov[p]\, w$.  Under spectral concentration,
$\Cov[p]$ has eigenvalues that are large only in the directions
corresponding to~$K_0$; aligning~$w$ with these directions while
suppressing the noisy tail reduces $\Var[s^{w}]$.  A lower-variance
score distribution yields a tighter gap between its median and its
$(1-\alpha)$-quantile, so $\qhat(w^{\star}) \leq \qhat(w^{1})$, and the
prediction band $\{y : \sSpec^{w^{\star}} \leq \qhat(w^{\star})\}$ is
correspondingly tighter.  The coverage guarantee is preserved because
the weights are fixed before calibration
(Section~\ref{sec:method:learned}).
\end{proof}

\begin{remark}
The three-way data split is essential for
Proposition~\ref{prop:width} to coexist with Theorem~\ref{thm:coverage}.
If the same data were used both to learn~$w$ and to compute~$\qhat$,
the calibration scores would no longer be exchangeable with the test
score, and the coverage guarantee would be invalidated.
\end{remark}


% ===========================================================================
\section{Method}
\label{sec:method}

% ---------------------------------------------------------------------------
\subsection{Problem Setup}
\label{sec:method:setup}

Let $\calX$ and $\calY$ denote spaces of functions defined on a bounded domain
$\Omega \subset \R^d$.  A neural operator $G_\theta : \calX \to \calY$,
parameterised by $\theta$, maps an input field (e.g., initial condition,
forcing term, or coefficient field) to a solution field of a partial
differential equation.  We train $G_\theta$ on a dataset
$\{(x_i, y_i)\}_{i=1}^{N}$ of input--solution pairs generated by a
numerical PDE solver.

After training, our goal is to construct \emph{prediction bands}
$\calC_\alpha(x) \subseteq \calY$ such that, for a new exchangeable test
point $(X_{n+1}, Y_{n+1})$,
\begin{equation}
  \label{eq:coverage}
  \Pr\!\bigl(Y_{n+1} \in \calC_\alpha(X_{n+1})\bigr) \;\geq\; 1 - \alpha,
\end{equation}
where $\alpha \in (0,1)$ is the user-specified miscoverage level.
Crucially, this guarantee must hold in finite samples and without
distributional assumptions on the data-generating process, beyond
exchangeability.

% ---------------------------------------------------------------------------
\subsection{Spectral Nonconformity Scores}
\label{sec:method:scores}

Given a prediction $\hat{y} = G_\theta(x)$ and the true solution $y$,
define the \emph{error field}
\begin{equation}
  e(x) \;=\; y - G_\theta(x).
\end{equation}
We analyse $e$ in the frequency domain via the discrete Fourier transform
(DFT):
\begin{equation}
  \hat{e}(k) \;=\; \mathrm{FFT}(e)(k), \qquad k \in \mathbb{Z}^d.
\end{equation}

\begin{definition}[Weighted Spectral Score]
\label{def:spectral-score}
For a weight function $w : \mathbb{Z}^d \to [0, \infty)$, the
\emph{weighted spectral nonconformity score} is
\begin{equation}
  \label{eq:spectral-score}
  \sSpec^{w}(G_\theta(x),\, y)
  \;=\; \sum_{k} w_k \,\bigl|\hat{e}(k)\bigr|^2.
\end{equation}
\end{definition}

Different choices of~$w$ recover well-known norms as special cases:
\begin{itemize}
  \item \textbf{$L^2$ score.}  Setting $w_k = 1$ for all~$k$ yields, by
    Parseval's theorem,
    $\sSpec^{1}(G_\theta(x), y) = \|y - G_\theta(x)\|_{L^2}^2$.
  \item \textbf{Sobolev $H^s$ score.}  Setting
    $w_k = (1 + |k|^2)^{s}$ for $s > 0$ yields the squared $H^s$
    Sobolev semi-norm of the error, which penalises high-frequency
    discrepancies more heavily.
  \item \textbf{Adaptive (learned) score.}  Setting $w_k$ to
    data-driven values (Section~\ref{sec:method:learned}) allows the
    score to concentrate on the frequency bands where the neural operator
    errs most, producing tighter prediction bands.
\end{itemize}

% ---------------------------------------------------------------------------
\subsection{Split Conformal Prediction with Spectral Scores}
\label{sec:method:conformal}

We apply the standard split conformal protocol
\citep{vovk2005algorithmic,lei2018distribution} with the spectral score
\eqref{eq:spectral-score} as the nonconformity measure.

Given a calibration set $\mathcal{D}_{\mathrm{cal}} = \{(x_i, y_i)\}_{i=1}^{n}$
that is exchangeable with the test point, we compute scores
$s_i = \sSpec^{w}(G_\theta(x_i), y_i)$ for $i = 1, \dots, n$ and set
\begin{equation}
  \label{eq:quantile}
  \qhat \;=\; \mathrm{Quantile}\!\Bigl(
    s_1, \dots, s_n;\;
    \frac{\lceil (1-\alpha)(n+1) \rceil}{n}
  \Bigr).
\end{equation}
The prediction band for a new input $x$ is then
\begin{equation}
  \label{eq:pred-band}
  \calC_\alpha(x)
  \;=\;
  \bigl\{\, y \in \calY : \sSpec^{w}(G_\theta(x),\, y) \leq \qhat \,\bigr\}.
\end{equation}
Because the spectral score is a measurable function of the residual, the
standard exchangeability-based coverage guarantee
(Theorem~\ref{thm:coverage}) applies immediately.

% ---------------------------------------------------------------------------
\subsection{Learned Spectral Weights}
\label{sec:method:learned}

The choice of weight vector~$w$ does not affect the \emph{validity} of
coverage (any measurable score preserves the guarantee), but it
strongly influences \emph{efficiency}: the average width of the
prediction band.  We therefore propose learning~$w$ from data.

\paragraph{Three-way data split.}
To avoid statistical double-dipping, we partition the held-out data into
three disjoint subsets:
\begin{enumerate}
  \item $\mathcal{D}_w$: used to learn the weight vector~$w$;
  \item $\mathcal{D}_{\mathrm{cal}}$: used to compute the conformal
    quantile $\qhat$ with the learned weights;
  \item $\mathcal{D}_{\mathrm{test}}$: used for evaluation only.
\end{enumerate}
The three-way split is necessary because using a single calibration set both
to optimise~$w$ and to compute~$\qhat$ would violate exchangeability
between the calibration scores and the test score, invalidating the
coverage guarantee.

\paragraph{Objective.}
On $\mathcal{D}_w$, we compute per-sample spectral error power spectra
and seek weights that minimise the \emph{variance} of the resulting scores:
\begin{equation}
  \label{eq:weight-obj}
  w^{\star}
  \;=\;
  \operatorname*{arg\,min}_{w \geq 0}
  \;\Var_{(x,y) \in \mathcal{D}_w}\!\bigl[\sSpec^{w}(G_\theta(x), y)\bigr].
\end{equation}
Minimising score variance concentrates the score distribution, which in turn
yields a smaller conformal quantile $\qhat$ and therefore tighter prediction
bands for a given coverage level~$1-\alpha$.

\paragraph{Binned parameterisation.}
To keep the optimisation low-dimensional and well-conditioned, we do not
learn one weight per frequency~$k$.  Instead, we partition the frequency
indices into $B = 16$ logarithmically spaced bins and assign a single
scalar weight to all frequencies within each bin:
$w_k = w_{b(k)}$ where $b(k) \in \{1, \dots, B\}$.
The resulting 16-dimensional optimisation is solved by L-BFGS-B with
non-negativity constraints.

% ---------------------------------------------------------------------------
\subsection{Algorithm}
\label{sec:method:algorithm}

Algorithm~\ref{alg:spectral-conformal} summarises the complete pipeline.

\begin{algorithm}[t]
\DontPrintSemicolon
\SetAlgoLined
\KwIn{Training data $\mathcal{D}_{\mathrm{train}}$; held-out data
  $\mathcal{D}_{\mathrm{held}}$; miscoverage level $\alpha$;
  number of frequency bins $B$}
\KwOut{Prediction band function $\calC_\alpha(\cdot)$}

\BlankLine
\tcp{Step 1: Train neural operator}
Train $G_\theta$ on $\mathcal{D}_{\mathrm{train}}$\;

\BlankLine
\tcp{Step 2: Three-way split of held-out data}
$\mathcal{D}_w,\, \mathcal{D}_{\mathrm{cal}},\, \mathcal{D}_{\mathrm{test}}
  \leftarrow \mathrm{Split}(\mathcal{D}_{\mathrm{held}})$\;

\BlankLine
\tcp{Step 3: Learn spectral weights}
\For{$(x_i, y_i) \in \mathcal{D}_w$}{
  $e_i \leftarrow y_i - G_\theta(x_i)$\;
  $\hat{e}_i \leftarrow \mathrm{FFT}(e_i)$\;
  Accumulate $|\hat{e}_i(k)|^2$ per frequency bin\;
}
$w^{\star} \leftarrow \operatorname*{arg\,min}_{w \geq 0}
  \Var\!\bigl[\sSpec^{w}\bigr]$ on $\mathcal{D}_w$
  \tcp*{Eq.~\eqref{eq:weight-obj}}

\BlankLine
\tcp{Step 4: Calibrate}
\For{$(x_i, y_i) \in \mathcal{D}_{\mathrm{cal}}$}{
  $s_i \leftarrow \sSpec^{w^{\star}}(G_\theta(x_i),\, y_i)$
    \tcp*{Eq.~\eqref{eq:spectral-score}}
}
$\qhat \leftarrow \mathrm{Quantile}\!\bigl(
  s_1, \dots, s_n;\;
  \lceil(1-\alpha)(n+1)\rceil / n\bigr)$
  \tcp*{Eq.~\eqref{eq:quantile}}

\BlankLine
\tcp{Step 5: Predict}
\For{new input $x$}{
  $\calC_\alpha(x) \leftarrow
    \{y : \sSpec^{w^{\star}}(G_\theta(x), y) \leq \qhat\}$\;
}
\Return{$\calC_\alpha$}\;

\caption{Spectral Conformal Prediction for Neural Operators}
\label{alg:spectral-conformal}
\end{algorithm}


% ===========================================================================
\section{Theoretical Analysis}
\label{sec:theory}

We establish three results: (i)~the spectral score with uniform weights
recovers the standard $L^2$ score, (ii)~the coverage guarantee of split
conformal prediction extends to all spectral scores, and (iii)~spectral
weighting can provably tighten prediction bands when errors exhibit
frequency concentration.

% ---------------------------------------------------------------------------
\begin{theorem}[Parseval Consistency]
\label{thm:parseval}
For uniform weights $w_k = 1$ for all~$k$, the weighted spectral score
equals the squared $L^2$ norm of the residual:
\begin{equation}
  \sSpec^{1}(G_\theta(x),\, y)
  \;=\;
  \|y - G_\theta(x)\|_{L^2(\Omega)}^2.
\end{equation}
\end{theorem}

\begin{proof}
By Parseval's theorem for the discrete Fourier transform,
\begin{equation}
  \sum_{k} \bigl|\hat{e}(k)\bigr|^2
  \;=\;
  \sum_{j} \bigl|e(j)\bigr|^2
  \;=\;
  \|e\|_{L^2}^2,
\end{equation}
where $e = y - G_\theta(x)$.  Substituting into
Definition~\ref{def:spectral-score} with $w_k = 1$ yields the result.
\end{proof}

% ---------------------------------------------------------------------------
\begin{theorem}[Coverage Guarantee]
\label{thm:coverage}
Let $s : \calX \times \calY \to \R$ be any measurable nonconformity score
function, and let $(X_1, Y_1), \dots, (X_n, Y_n), (X_{n+1}, Y_{n+1})$
be exchangeable.  Define $s_i = s(X_i, Y_i)$ and
\begin{equation}
  \qhat = \inf\!\Bigl\{q : \frac{|\{i \leq n : s_i \leq q\}|}{n}
    \geq \frac{\lceil (1-\alpha)(n+1) \rceil}{n}\Bigr\}.
\end{equation}
Then the prediction set $\calC_\alpha(x) = \{y : s(x,y) \leq \qhat\}$
satisfies
\begin{equation}
  \Pr\!\bigl(Y_{n+1} \in \calC_\alpha(X_{n+1})\bigr)
  \;\geq\;
  1 - \alpha.
\end{equation}
\end{theorem}

\begin{proof}
The proof follows directly from the standard split conformal guarantee
\citep{vovk2005algorithmic,lei2018distribution}.  The key observation is
that exchangeability of the data pairs $(X_i, Y_i)$ implies
exchangeability of the scores $s_i$, since $s$ is a fixed measurable
function applied identically to each pair.  In particular, the spectral
score $\sSpec^{w}$ is measurable because the FFT is a linear (hence
measurable) map and the weighted sum of squared moduli is a continuous
function of the Fourier coefficients.  The rank of $s_{n+1}$ among
$\{s_1, \dots, s_n, s_{n+1}\}$ is therefore uniformly distributed over
$\{1, \dots, n+1\}$, from which the coverage bound follows by the
standard quantile argument.
\end{proof}

% ---------------------------------------------------------------------------
\begin{proposition}[Band Width Reduction via Spectral Concentration]
\label{prop:width}
Let $\hat{e}_1, \dots, \hat{e}_n$ be the Fourier-domain error fields on
the calibration set.  Suppose that the spectral energy of the errors is
concentrated in a subset of frequency indices $K_0 \subset \mathbb{Z}^d$:
\begin{equation}
  \frac{\E\!\bigl[\sum_{k \in K_0} |\hat{e}(k)|^2\bigr]}
       {\E\!\bigl[\sum_{k} |\hat{e}(k)|^2\bigr]}
  \;\geq\; 1 - \delta,
  \qquad \delta \ll 1.
\end{equation}
Let $w^{\star}$ be the variance-minimising weights
\eqref{eq:weight-obj} and $w^{1}_k = 1$ the uniform weights.
Then the conformal quantile satisfies
$\qhat(w^{\star}) \leq \qhat(w^{1})$, with equality only when the
spectral energy is uniformly distributed across all frequencies.
Consequently, the prediction band under $w^{\star}$ is no wider than
under uniform weights, at the same coverage level~$1-\alpha$.
\end{proposition}

\begin{proof}[Proof sketch]
When errors concentrate in $K_0$, the spectral power outside $K_0$
contributes primarily noise variance to the scores under uniform weights.
By assigning larger weights to $K_0$ and down-weighting the remaining
frequencies, the optimised score $\sSpec^{w^{\star}}$ has strictly lower
variance than $\sSpec^{1}$ (the overall scale of scores is absorbed by the
quantile computation).

More precisely, write $s^{w} = w^{\!\top} p$ where
$p_k = |\hat{e}(k)|^2$ is the spectral power vector.  Then
$\Var[s^{w}] = w^{\!\top} \Cov[p]\, w$.  Under spectral concentration,
$\Cov[p]$ has eigenvalues that are large only in the directions
corresponding to~$K_0$; aligning~$w$ with these directions while
suppressing the noisy tail reduces $\Var[s^{w}]$.  A lower-variance
score distribution yields a tighter gap between its median and its
$(1-\alpha)$-quantile, so $\qhat(w^{\star}) \leq \qhat(w^{1})$, and the
prediction band $\{y : \sSpec^{w^{\star}} \leq \qhat(w^{\star})\}$ is
correspondingly tighter.  The coverage guarantee is preserved because
the weights are fixed before calibration
(Section~\ref{sec:method:learned}).
\end{proof}

\begin{remark}
The three-way data split is essential for
Proposition~\ref{prop:width} to coexist with Theorem~\ref{thm:coverage}.
If the same data were used both to learn~$w$ and to compute~$\qhat$,
the calibration scores would no longer be exchangeable with the test
score, and the coverage guarantee would be invalidated.
\end{remark}


% ===========================================================================
\section{Method}
\label{sec:method}

% ---------------------------------------------------------------------------
\subsection{Problem Setup}
\label{sec:method:setup}

Let $\calX$ and $\calY$ denote spaces of functions defined on a bounded domain
$\Omega \subset \R^d$.  A neural operator $G_\theta : \calX \to \calY$,
parameterised by $\theta$, maps an input field (e.g., initial condition,
forcing term, or coefficient field) to a solution field of a partial
differential equation.  We train $G_\theta$ on a dataset
$\{(x_i, y_i)\}_{i=1}^{N}$ of input--solution pairs generated by a
numerical PDE solver.

After training, our goal is to construct \emph{prediction bands}
$\calC_\alpha(x) \subseteq \calY$ such that, for a new exchangeable test
point $(X_{n+1}, Y_{n+1})$,
\begin{equation}
  \label{eq:coverage}
  \Pr\!\bigl(Y_{n+1} \in \calC_\alpha(X_{n+1})\bigr) \;\geq\; 1 - \alpha,
\end{equation}
where $\alpha \in (0,1)$ is the user-specified miscoverage level.
Crucially, this guarantee must hold in finite samples and without
distributional assumptions on the data-generating process, beyond
exchangeability.

% ---------------------------------------------------------------------------
\subsection{Spectral Nonconformity Scores}
\label{sec:method:scores}

Given a prediction $\hat{y} = G_\theta(x)$ and the true solution $y$,
define the \emph{error field}
\begin{equation}
  e(x) \;=\; y - G_\theta(x).
\end{equation}
We analyse $e$ in the frequency domain via the discrete Fourier transform
(DFT):
\begin{equation}
  \hat{e}(k) \;=\; \mathrm{FFT}(e)(k), \qquad k \in \mathbb{Z}^d.
\end{equation}

\begin{definition}[Weighted Spectral Score]
\label{def:spectral-score}
For a weight function $w : \mathbb{Z}^d \to [0, \infty)$, the
\emph{weighted spectral nonconformity score} is
\begin{equation}
  \label{eq:spectral-score}
  \sSpec^{w}(G_\theta(x),\, y)
  \;=\; \sum_{k} w_k \,\bigl|\hat{e}(k)\bigr|^2.
\end{equation}
\end{definition}

Different choices of~$w$ recover well-known norms as special cases:
\begin{itemize}
  \item \textbf{$L^2$ score.}  Setting $w_k = 1$ for all~$k$ yields, by
    Parseval's theorem,
    $\sSpec^{1}(G_\theta(x), y) = \|y - G_\theta(x)\|_{L^2}^2$.
  \item \textbf{Sobolev $H^s$ score.}  Setting
    $w_k = (1 + |k|^2)^{s}$ for $s > 0$ yields the squared $H^s$
    Sobolev semi-norm of the error, which penalises high-frequency
    discrepancies more heavily.
  \item \textbf{Adaptive (learned) score.}  Setting $w_k$ to
    data-driven values (Section~\ref{sec:method:learned}) allows the
    score to concentrate on the frequency bands where the neural operator
    errs most, producing tighter prediction bands.
\end{itemize}

% ---------------------------------------------------------------------------
\subsection{Split Conformal Prediction with Spectral Scores}
\label{sec:method:conformal}

We apply the standard split conformal protocol
\citep{vovk2005algorithmic,lei2018distribution} with the spectral score
\eqref{eq:spectral-score} as the nonconformity measure.

Given a calibration set $\mathcal{D}_{\mathrm{cal}} = \{(x_i, y_i)\}_{i=1}^{n}$
that is exchangeable with the test point, we compute scores
$s_i = \sSpec^{w}(G_\theta(x_i), y_i)$ for $i = 1, \dots, n$ and set
\begin{equation}
  \label{eq:quantile}
  \qhat \;=\; \mathrm{Quantile}\!\Bigl(
    s_1, \dots, s_n;\;
    \frac{\lceil (1-\alpha)(n+1) \rceil}{n}
  \Bigr).
\end{equation}
The prediction band for a new input $x$ is then
\begin{equation}
  \label{eq:pred-band}
  \calC_\alpha(x)
  \;=\;
  \bigl\{\, y \in \calY : \sSpec^{w}(G_\theta(x),\, y) \leq \qhat \,\bigr\}.
\end{equation}
Because the spectral score is a measurable function of the residual, the
standard exchangeability-based coverage guarantee
(Theorem~\ref{thm:coverage}) applies immediately.

% ---------------------------------------------------------------------------
\subsection{Learned Spectral Weights}
\label{sec:method:learned}

The choice of weight vector~$w$ does not affect the \emph{validity} of
coverage (any measurable score preserves the guarantee), but it
strongly influences \emph{efficiency}: the average width of the
prediction band.  We therefore propose learning~$w$ from data.

\paragraph{Three-way data split.}
To avoid statistical double-dipping, we partition the held-out data into
three disjoint subsets:
\begin{enumerate}
  \item $\mathcal{D}_w$: used to learn the weight vector~$w$;
  \item $\mathcal{D}_{\mathrm{cal}}$: used to compute the conformal
    quantile $\qhat$ with the learned weights;
  \item $\mathcal{D}_{\mathrm{test}}$: used for evaluation only.
\end{enumerate}
The three-way split is necessary because using a single calibration set both
to optimise~$w$ and to compute~$\qhat$ would violate exchangeability
between the calibration scores and the test score, invalidating the
coverage guarantee.

\paragraph{Objective.}
On $\mathcal{D}_w$, we compute per-sample spectral error power spectra
and seek weights that minimise the \emph{variance} of the resulting scores:
\begin{equation}
  \label{eq:weight-obj}
  w^{\star}
  \;=\;
  \operatorname*{arg\,min}_{w \geq 0}
  \;\Var_{(x,y) \in \mathcal{D}_w}\!\bigl[\sSpec^{w}(G_\theta(x), y)\bigr].
\end{equation}
Minimising score variance concentrates the score distribution, which in turn
yields a smaller conformal quantile $\qhat$ and therefore tighter prediction
bands for a given coverage level~$1-\alpha$.

\paragraph{Binned parameterisation.}
To keep the optimisation low-dimensional and well-conditioned, we do not
learn one weight per frequency~$k$.  Instead, we partition the frequency
indices into $B = 16$ logarithmically spaced bins and assign a single
scalar weight to all frequencies within each bin:
$w_k = w_{b(k)}$ where $b(k) \in \{1, \dots, B\}$.
The resulting 16-dimensional optimisation is solved by L-BFGS-B with
non-negativity constraints.

% ---------------------------------------------------------------------------
\subsection{Algorithm}
\label{sec:method:algorithm}

Algorithm~\ref{alg:spectral-conformal} summarises the complete pipeline.

\begin{algorithm}[t]
\DontPrintSemicolon
\SetAlgoLined
\KwIn{Training data $\mathcal{D}_{\mathrm{train}}$; held-out data
  $\mathcal{D}_{\mathrm{held}}$; miscoverage level $\alpha$;
  number of frequency bins $B$}
\KwOut{Prediction band function $\calC_\alpha(\cdot)$}

\BlankLine
\tcp{Step 1: Train neural operator}
Train $G_\theta$ on $\mathcal{D}_{\mathrm{train}}$\;

\BlankLine
\tcp{Step 2: Three-way split of held-out data}
$\mathcal{D}_w,\, \mathcal{D}_{\mathrm{cal}},\, \mathcal{D}_{\mathrm{test}}
  \leftarrow \mathrm{Split}(\mathcal{D}_{\mathrm{held}})$\;

\BlankLine
\tcp{Step 3: Learn spectral weights}
\For{$(x_i, y_i) \in \mathcal{D}_w$}{
  $e_i \leftarrow y_i - G_\theta(x_i)$\;
  $\hat{e}_i \leftarrow \mathrm{FFT}(e_i)$\;
  Accumulate $|\hat{e}_i(k)|^2$ per frequency bin\;
}
$w^{\star} \leftarrow \operatorname*{arg\,min}_{w \geq 0}
  \Var\!\bigl[\sSpec^{w}\bigr]$ on $\mathcal{D}_w$
  \tcp*{Eq.~\eqref{eq:weight-obj}}

\BlankLine
\tcp{Step 4: Calibrate}
\For{$(x_i, y_i) \in \mathcal{D}_{\mathrm{cal}}$}{
  $s_i \leftarrow \sSpec^{w^{\star}}(G_\theta(x_i),\, y_i)$
    \tcp*{Eq.~\eqref{eq:spectral-score}}
}
$\qhat \leftarrow \mathrm{Quantile}\!\bigl(
  s_1, \dots, s_n;\;
  \lceil(1-\alpha)(n+1)\rceil / n\bigr)$
  \tcp*{Eq.~\eqref{eq:quantile}}

\BlankLine
\tcp{Step 5: Predict}
\For{new input $x$}{
  $\calC_\alpha(x) \leftarrow
    \{y : \sSpec^{w^{\star}}(G_\theta(x), y) \leq \qhat\}$\;
}
\Return{$\calC_\alpha$}\;

\caption{Spectral Conformal Prediction for Neural Operators}
\label{alg:spectral-conformal}
\end{algorithm}


% ===========================================================================
\section{Theoretical Analysis}
\label{sec:theory}

We establish three results: (i)~the spectral score with uniform weights
recovers the standard $L^2$ score, (ii)~the coverage guarantee of split
conformal prediction extends to all spectral scores, and (iii)~spectral
weighting can provably tighten prediction bands when errors exhibit
frequency concentration.

% ---------------------------------------------------------------------------
\begin{theorem}[Parseval Consistency]
\label{thm:parseval}
For uniform weights $w_k = 1$ for all~$k$, the weighted spectral score
equals the squared $L^2$ norm of the residual:
\begin{equation}
  \sSpec^{1}(G_\theta(x),\, y)
  \;=\;
  \|y - G_\theta(x)\|_{L^2(\Omega)}^2.
\end{equation}
\end{theorem}

\begin{proof}
By Parseval's theorem for the discrete Fourier transform,
\begin{equation}
  \sum_{k} \bigl|\hat{e}(k)\bigr|^2
  \;=\;
  \sum_{j} \bigl|e(j)\bigr|^2
  \;=\;
  \|e\|_{L^2}^2,
\end{equation}
where $e = y - G_\theta(x)$.  Substituting into
Definition~\ref{def:spectral-score} with $w_k = 1$ yields the result.
\end{proof}

% ---------------------------------------------------------------------------
\begin{theorem}[Coverage Guarantee]
\label{thm:coverage}
Let $s : \calX \times \calY \to \R$ be any measurable nonconformity score
function, and let $(X_1, Y_1), \dots, (X_n, Y_n), (X_{n+1}, Y_{n+1})$
be exchangeable.  Define $s_i = s(X_i, Y_i)$ and
\begin{equation}
  \qhat = \inf\!\Bigl\{q : \frac{|\{i \leq n : s_i \leq q\}|}{n}
    \geq \frac{\lceil (1-\alpha)(n+1) \rceil}{n}\Bigr\}.
\end{equation}
Then the prediction set $\calC_\alpha(x) = \{y : s(x,y) \leq \qhat\}$
satisfies
\begin{equation}
  \Pr\!\bigl(Y_{n+1} \in \calC_\alpha(X_{n+1})\bigr)
  \;\geq\;
  1 - \alpha.
\end{equation}
\end{theorem}

\begin{proof}
The proof follows directly from the standard split conformal guarantee
\citep{vovk2005algorithmic,lei2018distribution}.  The key observation is
that exchangeability of the data pairs $(X_i, Y_i)$ implies
exchangeability of the scores $s_i$, since $s$ is a fixed measurable
function applied identically to each pair.  In particular, the spectral
score $\sSpec^{w}$ is measurable because the FFT is a linear (hence
measurable) map and the weighted sum of squared moduli is a continuous
function of the Fourier coefficients.  The rank of $s_{n+1}$ among
$\{s_1, \dots, s_n, s_{n+1}\}$ is therefore uniformly distributed over
$\{1, \dots, n+1\}$, from which the coverage bound follows by the
standard quantile argument.
\end{proof}

% ---------------------------------------------------------------------------
\begin{proposition}[Band Width Reduction via Spectral Concentration]
\label{prop:width}
Let $\hat{e}_1, \dots, \hat{e}_n$ be the Fourier-domain error fields on
the calibration set.  Suppose that the spectral energy of the errors is
concentrated in a subset of frequency indices $K_0 \subset \mathbb{Z}^d$:
\begin{equation}
  \frac{\E\!\bigl[\sum_{k \in K_0} |\hat{e}(k)|^2\bigr]}
       {\E\!\bigl[\sum_{k} |\hat{e}(k)|^2\bigr]}
  \;\geq\; 1 - \delta,
  \qquad \delta \ll 1.
\end{equation}
Let $w^{\star}$ be the variance-minimising weights
\eqref{eq:weight-obj} and $w^{1}_k = 1$ the uniform weights.
Then the conformal quantile satisfies
$\qhat(w^{\star}) \leq \qhat(w^{1})$, with equality only when the
spectral energy is uniformly distributed across all frequencies.
Consequently, the prediction band under $w^{\star}$ is no wider than
under uniform weights, at the same coverage level~$1-\alpha$.
\end{proposition}

\begin{proof}[Proof sketch]
When errors concentrate in $K_0$, the spectral power outside $K_0$
contributes primarily noise variance to the scores under uniform weights.
By assigning larger weights to $K_0$ and down-weighting the remaining
frequencies, the optimised score $\sSpec^{w^{\star}}$ has strictly lower
variance than $\sSpec^{1}$ (the overall scale of scores is absorbed by the
quantile computation).

More precisely, write $s^{w} = w^{\!\top} p$ where
$p_k = |\hat{e}(k)|^2$ is the spectral power vector.  Then
$\Var[s^{w}] = w^{\!\top} \Cov[p]\, w$.  Under spectral concentration,
$\Cov[p]$ has eigenvalues that are large only in the directions
corresponding to~$K_0$; aligning~$w$ with these directions while
suppressing the noisy tail reduces $\Var[s^{w}]$.  A lower-variance
score distribution yields a tighter gap between its median and its
$(1-\alpha)$-quantile, so $\qhat(w^{\star}) \leq \qhat(w^{1})$, and the
prediction band $\{y : \sSpec^{w^{\star}} \leq \qhat(w^{\star})\}$ is
correspondingly tighter.  The coverage guarantee is preserved because
the weights are fixed before calibration
(Section~\ref{sec:method:learned}).
\end{proof}

\begin{remark}
The three-way data split is essential for
Proposition~\ref{prop:width} to coexist with Theorem~\ref{thm:coverage}.
If the same data were used both to learn~$w$ and to compute~$\qhat$,
the calibration scores would no longer be exchangeable with the test
score, and the coverage guarantee would be invalidated.
\end{remark}
